\documentclass[12pt]{article}
\usepackage{../paper,../math}
\usepackage{../uark_colors}
% \addbibresource{references.bib}
\hypersetup{
  pdftitle = {UARK ECON 5783 AI Statement},
  pdfauthor = {Kyle Butts},
}

\hypersetup{
  colorlinks = true,
  allcolors = ozark_mountains,
  linkbordercolor = ozark_mountains,
  breaklinks = true,
  bookmarksopen = true
}

\usepackage{tcolorbox} 
\newtcolorbox{callout}[2][]{
  colback = apple_blossom,    % Background color
  coltext = ozark_mountains,  % Text color
  boxrule = 1pt,              % Border
  colframe = black!17,        % Border color 
  left = 1em,                 % Left padding
  right = 1em,                % Right padding
  top = 1em,                  % Top padding
  bottom = 0.75em,            % Bottom padding
  arc = 2mm                   % Rounded corners
  title = #2,                 % #2: title
  #1,                         % #1: tcolorbox options
}

% \usepackage{fontawesome} % For icons
\usepackage{setspace, titling}
\title{
  \vspace{-2em}
	{\huge \ttfamily Statement on AI} \\[-0.75em]
  {\Large \ttfamily [ECON 5783]} \\[-0.5em]
	{\Large Fall 2026}
}
\author{}
\date{}  % Toggle commenting to test

\begin{document}
\maketitle

This document details my thoughts on the use of large language models (`AI') in the classroom. 
Explicitly, I am choosing not to `ban' their use in the classroom for two reasons:
(i) much like calculators in the 1980s, these will be used in your career and you need to start learning to use them, and much more importantly, 
(ii) they are an incredibly helpful tool in \emph{the process of learning}.

Regularly, I use LLM when trying to learn something:
\begin{itemize}
  \item Ask questions: ``can you give me an introduction to X?''
  
  \item Clarify terminology: ``when they write `X', what do they mean?''
  
  \item Confirm my understanding: ``is this the right idea: ...?''
  
  \item Fix coding mistakes: ``This is my code and here is the error I am getting. Can you help me correct it?''
\end{itemize} 

It is really helpful! 
I can more quickly let my curiosity run wild and learn things I might have previously not had the energy/time to learn. 
From a liberal arts perspective, it is \emph{incredible} you have a personal tutor for free than can let you engage and learn from hundreds of years of scholarly thinking.
Because of the benefits to your learning, my job as your educator is to figure out how to get you these benefits without hurting your learning experiences.

\subsection*{AI use in this Course}

For the first goal (`learning to use LLMs'), I have not much concern. 
LLMs are pretty good at taking abstract requests and producing \emph{something}. 
Some might argue that prompting well is a skill that can be grown, but this does not seem to be true anymore.
However, very quickly `prompt engineering' is becoming a relic of earlier lower quality models. 
So learning to use LLMs is likely of little value in your career. 
Anyone can type requests into a LLM, so that is not how you will earn good salaries.

For the second goal (`using LLMs to learn'), I have more to say.  
There are two ways you can use LLMs in this class:
\begin{enumerate}
  \item \textbf{`Cheating model'} Hand off assignments to the LLM and it will produce something passable, perhaps even good.
  
  \item \textbf{`Tutor model'} Work with the LLM to complete the assignment in a way that you are able to learn more while completing the assignment. 
\end{enumerate}
There is a sort of continuum between these two extremes: somewhere in the middle, you are asking some questions, but ultimately just taking the solution given to you. 
Of course, the Cheating model is far less effort and given your competing interests, it becomes easy to slide along the continuum towards the Cheating model.

It is my job, therefore, to try and provide a counter-force that increases the value of the `Tutor model'. 
This semester, I am making my first attempt at this, but I need your help to gauge the effectiveness of this. 
The pitch: 
\begin{itemize}
  \item I will continue to give you take-home assignments that you will be required to turn in. 
  These will complement the lecture material and make sure you have a chance to work with real data and engage with the concepts in a hands-on way. 

  \item The class following the assignment's due date, we will have a short in-class quiz that will be based on the take-home assignment. 
\end{itemize}
The grade of the assignment will be a combination of the completed assignment and your grade on the in-class quiz. 
My hope is that the in-class quiz will require you to engage with the assignment more in the `Tutor model' and better learn from it.


\subsection*{Is it worth learning?}

A final thing worth discussing is the existential question: `why should I want to learn if `AI' knows everything already?'. 

My experience is that LLMs are a \emph{multiplier}. 
Production is something like this:
$$
  \text{Firm Productivity} = \text{Firm's Accumulated Judgement} \times \text{LLM Productivity Multiplier}
$$
If this equation is true, then it is obvious why you want to learn. 
Relative to before, learning is actually \emph{more valuable}, since you have the LLM multiplier to increase your abilities.

So is this equation true? 
For causal inference, I think this is absolutely true. 
The cost of running a regression and producing figures is near 0 now. 
There are 100s of models you can train, 100s of figures you can make, dashboards galore!! 

What's scarce isn't information. It's knowing which information to trust.

There's a famous book, How to Lie with Statistics, that shows how easy it is to use numbers to deceive people. 
LLMs create a closely related problem: it is now trivially easy to produce confident, well-formatted information that is completely wrong. 
And because it looks right, you believe it. 
You've \emph{lied with statistics to yourself}, and you might not even notice.

Think of it like this. 
An LLM is a research assistant who works incredibly fast, sounds authoritative, and produces valuable insights 20\% of the time. 
The hard part is not getting the output. 
The hard part is knowing which 20\% is useful.

One of the central goals of this class is learning to do exactly that: to read model output and figures critically. 
Learn to ask what's missing or misleading. 
That is the value-add I hope you will learn in this class.

More generally, I have a simple explanation for why I think that humans are essential in the production function.
Assume that companies like Google, OpenAI, and Anthropic want to make money (they, of course, want you to think they have more noble intentions). 
If their models casn do all that we as people can do, then what would be more profitable:
\begin{enumerate}
  \item Tell a LLM to start a company and make profits
  
  \item Sell `tokens' to companies and allow them to collect the excess profit. 
\end{enumerate}
Clearly, the former is more profitable as the excess value that a company makes would accrue to one of the big AI models rather than companies.
The reason they have to sell tokens and not just start a bunch of companies themselves is that without discernment, they are multiplying by 0 and end up with no productivity. 





\end{document}
